% ============================================================================
% Seccion 7: Evaluacion de la Aplicacion
% TFP - Rubrica: 2 puntos
% ============================================================================
\section{Evaluaci\'on de la aplicaci\'on}

La evaluaci\'on combina indicadores t\'ecnicos y de negocio. Dado que en esta
etapa el sistema es un prototipo, se describe el \emph{marco de evaluaci\'on}
definido; los valores medidos se completar\'an con las pruebas finales.

\subsection{M\'etricas empleadas}

\paragraph{Recuperaci\'on y ranking.} Para el motor de scoring (ranking Top-N) se
emplean m\'etricas como \textbf{Recall@K} y relevancia frente a un conjunto de
casos de referencia, evaluando si las carreras pertinentes aparecen en el Top-N.

\paragraph{Extractor de preferencias (LLM).} Se mide la correcci\'on de los
campos estructurados (regi\'on, presupuesto, pesos, c\'odigo RIASEC) y la tasa de
JSON malformado.

\paragraph{Justificaci\'on en lenguaje natural (LLM).} Se usa un componente de
\textbf{LLM-as-judge} que puntu\'a las explicaciones del agente, complementado con
evaluaci\'on humana de claridad, relevancia y \emph{factualidad} (que la
explicaci\'on se sustente en los datos del ranking y no invente cifras).

\paragraph{Negocio.} Satisfacci\'on del usuario, pertinencia percibida de la
recomendaci\'on y, a futuro, tasa de aceptaci\'on de la carrera sugerida.

\subsection{Costos como KPI}

Se consideran KPI de costo el costo mensual de infraestructura (con \emph{AWS
Budgets} como control) y el consumo de tokens del LLM (observado v\'ia LangSmith).
Estos indicadores son relevantes por tratarse de un proyecto acad\'emico con
presupuesto acotado.

\subsection{Casos de prueba y validaci\'on}

La evaluaci\'on se apoya en un dataset de casos representativos que cubren
distintos perfiles de estudiante (por regi\'on, presupuesto y prioridades), sobre
los cuales se ejecuta el flujo completo (extracci\'on de preferencias
$\rightarrow$ scoring $\rightarrow$ justificaci\'on) y se contrastan las salidas
con lo esperado. Se contemplan del orden de \textbf{5--6 casos can\'onicos} (por
ejemplo, el estudiante de Cusco con presupuesto acotado de la Secci\'on~8).

\paragraph{Reto de los \emph{evals} en un sistema conversacional.} A diferencia de
una interacci\'on simple de pregunta--respuesta, aqu\'i la conversaci\'on es
\textbf{multi-turno}: para que un \emph{eval} sea representativo debe simular
varios mensajes del estudiante (saludo, aporte de preferencias, aclaraciones),
no un \'unico JSON de entrada. Por ello los casos de prueba se construyen como
\emph{guiones de conversaci\'on} y se eval\'uan con el \emph{LLM-as-judge}. Este
componente se implementa del lado de AWS y queda \textbf{pendiente} hasta que el
agente est\'e desplegado.

\subsection{Comparaci\'on de enfoques (resultados)}

Se compar\'o el enfoque propuesto (pesos din\'amicos inferidos por el LLM) frente a
dos \emph{baselines}, midiendo cu\'antas carreras del Top-5 comparte cada uno con el
enfoque propuesto para el caso de Cusco (288 candidatas tras los filtros):

\begin{table}[H]
\centering
\begin{tabular}{@{}lcl@{}}
\toprule
Enfoque & Top-5 en com\'un & Interpretaci\'on \\
\midrule
Pesos din\'amicos (propuesto) & --- & referencia \\
Pesos iguales (baseline 1)    & 1/5 & el ranking cambia casi por completo \\
Solo ingreso (baseline 2)     & 3/5 & los pesos s\'i modifican el resultado \\
\bottomrule
\end{tabular}
\caption{El motor de 5 factores no equivale a un \emph{baseline} simple.}
\label{tab:comparacion-enfoques}
\end{table}

\paragraph{Prueba de consistencia.} Para dos perfiles en Lima, uno que prioriza el
costo y otro el ingreso, el \'indice de accesibilidad de costo (\texttt{cost\_norm},
mayor $=$ m\'as accesible) del Top-10 fue \textbf{0.941} vs.\ \textbf{0.502}: el
motor entrega opciones m\'as accesibles a quien prioriza el costo, confirmando que
el ranking responde a las preferencias del estudiante.

\paragraph{Nota metodol\'ogica.} Estos resultados eval\'uan el \emph{motor de
recomendaci\'on} (determin\'istico) sobre \texttt{features.csv}; la evaluaci\'on de
las respuestas en lenguaje natural del agente (\emph{LLM-as-judge},
\emph{groundedness}) se realiza sobre el sistema desplegado.

\subsection{Limitaciones detectadas}

\begin{itemize}[noitemsep]
  \item Los datos de salario reflejan solo el empleo formal y pueden subestimar
        ingresos reales en ciertos rubros.
  \item Las carreras o universidades con pocos egresados producen promedios poco
        confiables.
  \item Los pesos y la f\'ormula de scoring de esta etapa son ilustrativos y
        requieren calibraci\'on con usuarios reales.
  \item El LLM puede generar pesos incoherentes ante mensajes ambiguos; por ello
        la explicaci\'on se ancla siempre en datos verificables del ranking.
  \item La evaluaci\'on humana est\'a limitada a un n\'umero reducido de usuarios
        de prueba en esta fase.
  \item La variable de \textbf{costo anual} es la menos confiable del dataset:
        aproximadamente el \textbf{50\%} de los registros est\'a vac\'io o imputado
        (19\% nulos, 31\% en cero, 50\% imputado). El motor usa la versi\'on
        normalizada e imputada (\texttt{cost\_norm}), por lo que el \emph{ranking}
        se mantiene v\'alido, pero los montos absolutos de costo no deben
        interpretarse de forma literal.
\end{itemize}
